\documentclass[11pt]{report}
% \documentclass[12pt]{amsart}
% \documentclass[12pt]{amsproc}
% \documentclass[12pt]{amsbook}

\usepackage{amssymb,amscd,amsthm, verbatim,amsmath,color,fancyhdr, mathrsfs}
\usepackage{graphicx}
\usepackage{turnstile}
\usepackage[hidelinks]{hyperref}

\def\tequiv{\ensuremath\sdststile{}{}}

\usepackage[letterpaper, left=2cm, right=2cm, top=2.5cm,
bottom=2.5cm,dvips]{geometry}

\setcounter{section}{1}

\newtheorem{thm}{Theorem}[section]
\newtheorem{prop}[thm]{Proposition}
\newtheorem{lem}[thm]{Lemma}
\newtheorem{rem}[thm]{Remark}
\newtheorem{cor}[thm]{Corollary}
\newtheorem{exr}[thm]{Exercise}
\newtheorem{ex}[thm]{Example}
\newtheorem{defn}[thm]{Definition}

\newcommand{\R}{\mathbb{R}}
\newcommand{\Z}{\mathbb{Z}}
\newcommand{\N}{\mathbb{N}}
\newcommand{\aB}{\mathcal{B}}
\newcommand{\aM}{\mathcal{M}}

\title{Measure and Probability theory}

\author{Victor Schmit}

\date{\today}

\begin{document}
\maketitle

\tableofcontents
\newpage

\chapter{Measures}

\section{Lebesgue Measure}

There are several ways to define the Lebesgue measure, denoted by $\lambda$.
We propose a construction based on outer measures for the one–dimensional case.
In a first part, we construct the measure on $\mathbb{R}$ and then extend the
result to the multidimensional case using the classical outer measure definition.

There is a more natural way to extend the result to higher dimensions using
product measures. However, the approach presented here provides an opportunity
to examine some proofs that are rarely discussed when dealing directly with
multidimensional constructions.

\subsection{The one–dimensional case}

\defn Let us define
\[
	\lambda^*(A) =
	\inf\left\{
	\sum_{n=0}^{+\infty} (b_n-a_n)
	\ :\
	A \subset \bigcup_{n=0}^{+\infty} ]a_n,b_n[
	\right\}
\]
for any $A \subset \mathbb{R}$.

\prop The following properties hold:
\begin{enumerate}
	\item $\lambda^*$ is an outer measure.

	\item $\mathcal{B}(\mathbb{R}) \subset \mathcal{M}(\lambda^*)$.

	\item $\lambda^*(]a,b[)=b-a$.
\end{enumerate}

\rem Although the definition of $\lambda^*$ is given in terms
of countable coverings, we may of course also use finite coverings when
convenient.

Indeed, if a set $A$ is covered by finitely many intervals
\[
	A \subset \bigcup_{n=0}^N ]a_n,b_n[,
\]
then for any $\varepsilon>0$ we can complete this family into a countable
covering by adding intervals of arbitrarily small total length, for instance
$]0,\varepsilon 2^{-n}[$ for $n\ge N+1$. In this way we obtain a countable
covering of $A$ whose total length is
\[
	\sum_{n=0}^N (b_n-a_n) + \sum_{n=N+1}^{+\infty} \varepsilon 2^{-n}
	\le
	\sum_{n=0}^N (b_n-a_n)+\varepsilon.
\]

Since $\varepsilon>0$ is arbitrary, any estimate obtained from a finite
covering may be used in the definition of $\lambda^*$.

\begin{proof}
	Let $A$ be a subset of $\mathbb{R}$.

	\emph{(1)} Let us show that $\lambda^*$ is an outer measure.
	First, let us notice that $\lambda^*(\varnothing)=0$. Indeed,
	\[
	\varnothing \subset ]0,\varepsilon [.
		\]
		So,
		\[
			\lambda^*(\varnothing) \le \varepsilon.
		\]
		It is then clear that $\lambda^*(\varnothing) = 0$.
		Next, by basic properties of the infimum, if $A \subset B$ then
		\[
			\lambda^*(A)\le \lambda^*(B).
		\]
		It remains to prove countable subadditivity.
		Let $(A_n)_{n\in\mathbb{N}}$ be a sequence of subsets of $\mathbb{R}$.
		If
		\[
			\sum_{n=0}^{+\infty}\lambda^*(A_n)=+\infty,
		\]
		then clearly
		\[
			\lambda^*\!\left(\bigcup_{n=0}^{+\infty}A_n\right)
			\le
			\sum_{n=0}^{+\infty}\lambda^*(A_n).
		\]
		Assume now that the series is finite, so that
		$\lambda^*(A_n)<+\infty$ for every $n$.
		Fix $\varepsilon>0$.
		For each $n\in\mathbb{N}$ there exists a sequence of intervals
		$]a_{i,n},b_{i,n}[$ covering $A_n$ such that
		\[
			\sum_{i=0}^{+\infty}(b_{i,n}-a_{i,n})
			\le
			\lambda^*(A_n)+\varepsilon 2^{-n}.
		\]
		The family of intervals $(]a_{i,n},b_{i,n}[)_{i,n\in\mathbb{N}}$
		covers $\displaystyle\bigcup_{n=0}^{+\infty}A_n$.
		Therefore
		\[
			\lambda^*\!\left(\bigcup_{n=0}^{+\infty}A_n\right)
			\le
			\sum_{n=0}^{+\infty}\sum_{i=0}^{+\infty}(b_{i,n}-a_{i,n})
			\le
			\sum_{n=0}^{+\infty}\lambda^*(A_n)+\varepsilon.
		\]
		Since $\varepsilon>0$ is arbitrary, we obtain
		\[
			\lambda^*\!\left(\bigcup_{n=0}^{+\infty}A_n\right)
			\le
			\sum_{n=0}^{+\infty}\lambda^*(A_n).
		\]
		This proves countable subadditivity and therefore $\lambda^*$ is an outer measure.

		\medskip

		\emph{(2)} Let us show that $\mathcal{B}(\mathbb{R}) \subset \mathcal{M}(\lambda^*)$. Recall that $\mathcal{B}(\mathbb{R})=\sigma(]-\infty,a]\,:\,a\in\mathbb{R})$.
		It is therefore sufficient to show that every interval of the form
	$]-\infty,a]$ is $\aM(\lambda^*)$–measurable.
		Let $a\in\mathbb{R}$ and let $(]a_n,b_n[)_{n\in\mathbb{N}}$ be a countable
		family of intervals covering a set $A\subset\mathbb{R}$.
		Then we may decompose the sum of their lengths as
		\begin{align}
			\sum_{n=0}^{+\infty}(b_n-a_n)
			 & =
			\sum_{\substack{n\ge0 \\ b_n\le a}}(b_n-a_n)
			+
			\sum_{\substack{n\ge0 \\ a_n<a<b_n}}(b_n-a_n)
			+
			\sum_{\substack{n\ge0 \\ a\le a_n}}(b_n-a_n).
		\end{align}
		Let $\varepsilon>0$ and define
		\[
			c_n=a-\varepsilon2^{-n}, \qquad d_n=a+\varepsilon2^{-n}.
		\]
		Then
		\[
			4\varepsilon
			+
			\sum_{\substack{n\ge0 \\ a_n<a<b_n}}(b_n-a_n)
			\ge
			\sum_{\substack{n\ge0 \\ a_n<a<b_n}}(b_n-c_n)
			+
			\sum_{\substack{n\ge0 \\ a_n<a<b_n}}(d_n-a_n).
		\]
		Combining these inequalities yields
		\[
			4\varepsilon
			+
			\sum_{n=0}^{+\infty}(b_n-a_n)
			\ge
			\lambda^*(A\cap]-\infty,a])
			+
			\lambda^*(A\cap]a,+\infty[).
		\]
		Taking the infimum over all coverings of $A$ gives
		\[
			4\varepsilon+\lambda^*(A)
			\ge
			\lambda^*(A\cap]-\infty,a])
			+
			\lambda^*(A\cap]a,+\infty[).
		\]
		Since $\varepsilon>0$ is arbitrary, we obtain
		\[
			\lambda^*(A)
			\ge
			\lambda^*(A\cap]-\infty,a])
			+
			\lambda^*(A\cap]a,+\infty[).
		\]
		Hence $]-\infty,a]$ is $\aM(\lambda^*)$–measurable. Since the Borel
	$\sigma$–algebra is generated by such intervals, we conclude that
		\[
			\mathcal{B}(\mathbb{R})\subset\mathcal{M}(\lambda^*).
		\]

		\medskip

		\emph{(3)} Let us prove that $\lambda^*(]a,b[]) = b-a$.
		The interval $]a,b[$ itself covers $]a,b[$, then
		\[
			\lambda^*(]a,b[)\le b-a.
		\]
		To prove the other inequality, choose $a<c<d<b$. We have
		\[
			\lambda^*([c,d])\le \lambda^*(]a,b[).
		\]
		Let $(]a_n,b_n[))_{n\in\mathbb{N}}$ be a covering of $]a,b[$. Then it is also a covering of $[c,d]$, which is compact. Hence there exists $N\in\mathbb{N}$ such that the finite family $(]a_n,b_n[))_{0\le n\le N}$ covers $[c,d]$.
		Now let us denote
		\[
			D_m=\left\{L_k^m=[k2^{-m},(k+1)2^{-m}[\,:\,k\in\mathbb{Z}\right\}.
		\]
		Note that $|L_k^m|$ denotes the length of the interval $L_k^m$, hence
	$|L_k^m|=2^{-m}$.
		For any $n\in\{0,\dots,N\}$, let $A_{n,m}$ be the set of indices $k$
		such that
		\[
			L_k^m\cap ]a_n,b_n[\neq\varnothing .
		\]
		Observe that the intervals $(L_k^m)_{k\in A_{n,m}}$ are contiguous.
		Indeed, since $]a_n,b_n[$ is an interval, if $L_{k_1}^m$ and
	$L_{k_2}^m$ intersect $]a_n,b_n[$ with $k_1<k_2$, then every
	$L_k^m$ with $k_1\le k\le k_2$ also intersects $]a_n,b_n[$.
		Thus the union of these intervals is itself an interval.
		Moreover,
		\[
			]a_n,b_n[
			\subset
			\bigcup_{k\in A_{n,m}} L_k^m .
		\]
		This union may extend slightly outside $]a_n,b_n[$ near the two
		endpoints $a_n$ and $b_n$. At most two dyadic intervals may cross these
		boundary points, and each has length $2^{-m}$. Therefore
		\[
			\sum_{k\in A_{n,m}} |L_k^m|
			\le
			(b_n-a_n)+2\cdot 2^{-m}.
		\]
		Rearranging gives
		\[
			b_n-a_n
			\ge
			\sum_{k\in A_{n,m}} |L_k^m| - 2\cdot 2^{-m}.
		\]
		Summing over $n=0,\dots,N$ yields
		\[
			\sum_{n=0}^{N}(b_n-a_n)
			\ge
			\sum_{n=0}^{N}\sum_{k\in A_{n,m}} |L_k^m| - N2^{-m+1}.
		\]
		It is clear that the chosen intervals $L_k^m$ cover the union of the
		intervals $]a_n,b_n[$, hence they also cover $[c,d]$.
		Let us denote by $B_m$ the set of indices $k$ corresponding to the chosen
		intervals $L_k^m$. Then
		\[
			[c,d] \subset \bigcup_{k\in B_m} L_k^m
			= \bigcup_{n=0}^{N} \bigcup_{k \in A_{n,m}} L_k^m .
		\]
		From this family we can extract a subset $C_m\subset B_m$ such that the
		intervals $(L_k^m)_{k\in C_m}$ are contiguous and still cover $[c,d]$.
		Since $|L_k^m|$ denotes the length of $L_k^m$, and the intervals
	$(L_k^m)_{k\in C_m}$ form a contiguous family covering $[c,d]$, we have
		\[
			\sum_{k\in C_m} |L_k^m| \ge d-c.
		\]
		Moreover,
		\[
			\sum_{n=0}^{N}\sum_{k\in A_{n,m}} |L_k^m|
			\ge
			\sum_{k\in B_m} |L_k^m|
			\ge
			\sum_{k\in C_m} |L_k^m|.
		\]
		Therefore
		\[
			\sum_{n=0}^{N}\sum_{k\in A_{n,m}} |L_k^m|
			\ge d-c.
		\]
		Combining the previous inequalities, we obtain
		\[
			\sum_{n=0}^{N}(b_n-a_n)\ge d-c-N2^{-m+1}.
		\]
		Letting $m\to+\infty$, we get
		\[
			\sum_{n=0}^{N}(b_n-a_n)\ge d-c.
		\]
		Since
		\[
			\sum_{n=0}^{+\infty}(b_n-a_n)\ge \sum_{n=0}^{N}(b_n-a_n),
		\]
		it follows that every covering $(]a_n,b_n[)_{n\in\mathbb{N}}$ of $]a,b[$ satisfies
		\[
			\sum_{n=0}^{+\infty}(b_n-a_n)\ge d-c.
		\]
		Taking the infimum over all such coverings yields
		\[
			\lambda^*(]a,b[)\ge d-c.
		\]
		Since this holds for every $a<c<d<b$, we conclude by letting $c\to a$ and $d\to b$ that
	\[
		\lambda^*(]a,b[)\ge b-a.
	\]
	Together with the reverse inequality, this proves that
	\[
		\lambda^*(]a,b[)=b-a.
	\]
\end{proof}

\subsection{The multidimensional case}


\section{Stieljes Integral}
\label{sec:stieljes}

\section{Bounded Variation}

\defn[Total Variation]
Let $a : [0, T] \to \R$ be a function. For any $t \in [0,T]$, the \emph{total variation}
of $a$ on $[0,t]$ is defined by
\[
	V(t) =
	\sup \left\{
	\sum_{i=0}^{n}
	|a(t_{i+1}) - a(t_i)|
	\; : \;
	0 = t_0 < t_1 < \dots < t_{n+1} = t
	\right\}.
\]

\defn[Bounded Variation]
Let $a : [0, T] \to \R$ be a continuous function such that $a(0)=0$.
We say that $a$ is of \emph{bounded variation} on $[0,T]$ if
\[
	V(T) < +\infty .
\]

\thm
\label{thm:bound-stiel}
Let $a : [0,T] \to \R$ be continuous with $a(0)=0$.
Then $a$ is of bounded variation if and only if there exists a finite signed
measure $\mu$ on $[0,T]$ such that
\[
	a(t) = \mu([0,t]), \qquad t \in [0,T].
\]
In this case the total variation $V(T)$ coincides with the total variation of the
measure $\mu$.

\rem
This theorem is important because it allows us to interpret a function of bounded
variation as the cumulative mass of a signed measure. In other words, the
increments of the function correspond to the mass assigned by the measure:
\[
	a(t) - a(s) = \mu((s,t]).
\]

This provides a natural way to define integrals with respect to $a$.
If $f$ is a suitable function (for instance continuous), we define
\[
	\int_0^T f(t)\, da(t)
	:= \int_{[0,T]} f(t)\, \mu(dt).
\]

\rem
The integral with respect to a function of bounded variation can also be
defined as the limit of Riemann--Stieltjes sums.

Let
\[
	0 = t_0 < t_1 < \dots < t_n = T
\]
be a partition of $[0,T]$, and denote its mesh by
\[
	\max_{0 \le i \le n-1} (t_{i+1}-t_i).
\]

If $f$ is continuous and $a$ is of bounded variation, we define
\[
	\int_0^T f(t)\, da(t)
	=
	\lim_{\max (t_{i+1}-t_i) \to 0}
	\sum_{i=0}^{n-1}
	f(t_i)\big(a(t_{i+1})-a(t_i)\big).
\]

The bounded variation property ensures that the total size of the increments
\[
	\sum_{i=0}^{n-1} |a(t_{i+1})-a(t_i)|
\]
remains bounded. This prevents the Riemann sums from oscillating wildly and
ensures that the limit defining the integral exists.

To prove Theorem \ref{thm:bound-stiel}, we rely on several lemmas and
properties of the total variation function $V$.

\prop
Let $a : [0,T] \to \mathbb{R}$ and let $V$ denote its total variation
function.

\begin{enumerate}

	\item The function $V$ is non-decreasing and non-negative.
	\item For every $0 \le s < t \le T$,
	      \[
		      V(t) - V(s)
		      =
		      \sup \left\{
		      \sum_{i=0}^{n}
		      |a(t_{i+1}) - a(t_i)|
		      \ : \
		      s = t_0 < t_1 < \dots < t_{n+1} = t
		      \right\}.
	      \]

	      In other words, $V(t)-V(s)$ is the total variation of $a$ over the
	      interval $[s,t]$. We denote the total variation of $a$ over the interval
	      $[s,t]$ by $V_a(s,t)$.

	\item The function $V$ is continuous on $[0,T]$.
\end{enumerate}

\begin{proof}
	Let $a:[0,T]\to\R$ be a continuous function.

	\emph{(1)} We do not prove this point as it is very clear that $V$ is non-decreasing and non-negative.

	\medskip

	\emph{(2)} $V(t)-V(s)$ equals the variation of $a$ on $[s,t]$.
	Fix $0\le s < t\le T$ and define the variation of $a$ on $[s,t]$ by
	\[
		V_a(s,t)
		:=
		\sup\left\{
		\sum_{j=0}^{m-1}\bigl|a(x_{j+1})-a(x_j)\bigr|
		:\;
		s=x_0<\cdots<x_m=t
		\right\}.
	\]
	We prove that $V(t)-V(s)=V_a(s,t)$.

	\medskip

	\emph{Step 1: $V(t)-V(s)\ge V_a(s,t)$.}
	Let $\varepsilon>0$ and let $s=x_0<\cdots<x_m=t$ be an arbitrary partition of $[s,t]$.
	By definition of $V(s)$, there exists a partition $0=u_0<\cdots<u_n=s$ of $[0,s]$
	such that
	\[
		V(s)\le \sum_{i=0}^{n-1}\bigl|a(u_{i+1})-a(u_i)\bigr|+\varepsilon.
	\]
	Concatenate the two partitions to obtain a partition of $[0,t]$:
	\[
		0=u_0<\cdots<u_n=s=x_0<\cdots<x_m=t.
	\]
	Therefore, by definition of $V(t)$,
	\[
		V(t)\ge
		\sum_{i=0}^{n-1}\bigl|a(u_{i+1})-a(u_i)\bigr|
		+
		\sum_{j=0}^{m-1}\bigl|a(x_{j+1})-a(x_j)\bigr|.
	\]
	Subtracting the previous inequality for $V(s)$ yields
	\[
		V(t)-V(s)\ge
		\sum_{j=0}^{m-1}\bigl|a(x_{j+1})-a(x_j)\bigr|-\varepsilon.
	\]
	Since the partition of $[s,t]$ was arbitrary, taking the supremum over all such
	partitions gives
	\[
		V(t)-V(s)\ge V_a(s,t)-\varepsilon.
	\]
	Letting $\varepsilon\downarrow0$ we obtain $V(t)-V(s)\ge V_a(s,t)$.

	\medskip

	\emph{Step 2: $V(t)-V(s)\le V_a(s,t)$.}
	Let $\varepsilon>0$. By definition of $V(t)$, there exists a partition
	$0=t_0<\cdots<t_N=t$ of $[0,t]$ such that
	\[
		V(t)\le \sum_{k=0}^{N-1}\bigl|a(t_{k+1})-a(t_k)\bigr|+\varepsilon.
	\]
	Insert the point $s$ into this partition (if it is not already present), obtaining
	a refined partition $0=r_0<\cdots<r_M=t$ with $s=r_\ell$ for some $\ell$.
	Refinement can only increase the sum of increments, hence
	\[
		\sum_{k=0}^{N-1}\bigl|a(t_{k+1})-a(t_k)\bigr|
		\le
		\sum_{j=0}^{M-1}\bigl|a(r_{j+1})-a(r_j)\bigr|.
	\]
	Consequently,
	\[
		V(t)\le \sum_{j=0}^{M-1}\bigl|a(r_{j+1})-a(r_j)\bigr|+\varepsilon.
	\]
	Split the latter sum at $s=r_\ell$:
	\[
		\sum_{j=0}^{M-1}\bigl|a(r_{j+1})-a(r_j)\bigr|
		=
		\sum_{j=0}^{\ell-1}\bigl|a(r_{j+1})-a(r_j)\bigr|
		+
		\sum_{j=\ell}^{M-1}\bigl|a(r_{j+1})-a(r_j)\bigr|.
	\]
	The first sum is bounded by $V(s)$, and the second sum is bounded by $V_a(s,t)$.
	Therefore,
	\[
		V(t)\le V(s)+V_a(s,t)+\varepsilon,
	\]
	hence
	\[
		V(t)-V(s)\le V_a(s,t)+\varepsilon.
	\]
	Letting $\varepsilon\downarrow0$ gives $V(t)-V(s)\le V_a(s,t)$.
	Combining Steps 1 and 2 yields $V(t)-V(s)=V_a(s,t)$.

	\medskip

	\emph{(3)} $V$ is continuous. Let $0 \le s < t \le T$. Since $V$ is non–decreasing, the right limit
	\[
		V(s^+) := \lim_{u \downarrow s} V(u)
	\]
	and the left limit
	\[
		V(t^-) := \lim_{u \uparrow t} V(u)
	\]
	exist.
	Moreover,
	\[
		V(s) \le V(s^+) \le V(t^-) \le V(t).
	\]
	In particular,
	\[
		V(t^-) - V(s^+) \le V(t) - V(s).
	\]
	If we prove the reverse inequality, we obtain
	\[
		V(t)-V(s)=V(t^-)-V(s^+).
	\]
	Since $V(s)\le V(s^+)\le V(t^-)\le V(t)$, this would imply
	\[
		V(s^+)=V(s) \quad \text{and} \quad V(t^-)=V(t),
	\]
	hence $V$ would be continuous.
	We now prove the reverse inequality.
	Let $\varepsilon>0$ and let
	\[
		s=x_0 < x_1 < \dots < x_n=t
	\]
	be a partition of $[s,t]$.
	Since $a$ is continuous on $t$ and $s$, there exists $\delta>0$ such that
	\[
		|t-u|<\delta \quad \Longrightarrow \quad |a(t)-a(u)|<\varepsilon
		\quad \text{ and }\quad
		|s-v|<\delta \quad \Longrightarrow \quad |a(s)-a(v)|<\varepsilon .
	\]
	Refining the partition if necessary, we may assume that its mesh satisfies
	\[
		x_{k+1}-x_k<\delta .
	\]
	Now remove the endpoints $x_0=s$ and $x_n=t$ and denote
	\[
		y_0=x_1,\quad y_1=x_2,\quad \dots,\quad y_{n-2}=x_{n-1}.
	\]
	Then $(y_i)$ is a partition of $[y_0,y_{n-2}] \subset (s,t)$.
	Using the triangle inequality and the choice of $\delta$, we obtain
	\[
		\sum_{k=0}^{n-1}|a(x_{k+1})-a(x_k)|
		\le
		2\varepsilon + \sum_{k=0}^{n-3}|a(y_{k+1})-a(y_k)|.
	\]
	By the result of Point (2),
	\[
		\sum_{k=0}^{n-3}|a(y_{k+1})-a(y_k)|
		\le V(y_{n-2})-V(y_0).
	\]
	Since $y_0>s$ and $y_{n-2}<t$, we have
	\[
		V(y_{n-2})-V(y_0)\le V(t^-)-V(s^+).
	\]
	Therefore,
	\[
		\sum_{k=0}^{n-1}|a(x_{k+1})-a(x_k)|
		\le
		2\varepsilon + V(t^-)-V(s^+).
	\]
	Taking the supremum over all partitions $(x_i)$ of $[s,t]$ yields
	\[
		V(t)-V(s)\le 2\varepsilon + V(t^-)-V(s^+).
	\]
	Since $\varepsilon>0$ is arbitrary, we conclude that
	\[
		V(t)-V(s)\le V(t^-)-V(s^+).
	\]
	Combining this inequality with the reverse inequality obtained earlier gives
	\[
		V(t)-V(s)=V(t^-)-V(s^+).
	\]
	Hence $V$ is continuous on $[0,T]$.

\end{proof}

\rem
\label{rem:cont}
One can notice in the proof related to the continuity of the function $V$, that
$a$ is continuous to the right (to the left) implies that $V$ is continuous to the right (to the left).

\lem[Jordan decomposition]
\label{lem:jordan-cont}
Let $a:[0,T]\to\R$ be a continuous function of bounded variation with $a(0)=0$.
Then there exist two continuous non-decreasing functions $N,D:[0,T]\to\R$
such that
\[
	a(t)=N(t)-D(t), \qquad t\in[0,T],
\]
and
\[
	N(0)=D(0)=0.
\]

\begin{proof}
	Let $t \in [0,T]$. We define
	\[
		N(t) = \frac{V(t) + a(t)}{2}, \quad D(t) = \frac{V(t) - a(t)}{2}.
	\]
	Based on the definition of $V$, it is clear that $N,D$ are non-decreasing and non-negative.
	In addition, because $V$ and $a$ are continuous, $D$ and $N$ are also continuous. Finally,
	\[
		a(t)=N(t)-D(t)
	\]
	and
	\[
		N(0)=D(0)=0.
	\]
\end{proof}

\rem One can also notice that $V = N + D$.

\begin{proof}[Proof of Theorem \ref{thm:bound-stiel}]
	We want to prove that there exists a signed measure $\mu$ on $[0,T]$ such that $a(t) = \mu\left([0,t]\right)$.
	Based on the Jordan decomposition (Lemma \ref{lem:jordan-cont}), we have the following formula :
	\[
		a(t)=N(t)-D(t), \qquad t\in[0,T],
	\]
	with $N$ and $D$ continuous, non-decreasing, non-negative and that starts from $0$.
	Based on Section \ref{sec:stieljes} which deals with Stieljes integrals, there exists $\mu_1$ and $\mu_2$
	two positive measures on $[0,T]$ such that for every $t\in [0,T]$,
	\[
		N(t) = \mu_1([0,t]), \quad
		D(t) = \mu_2([0,t]).
	\]
	Finally, defining $\mu = \mu_1 - \mu_2$, a signed measure on $[0,T]$, we have
	\[
		a(t)=\mu([0,t]), \qquad t\in[0,T].
	\]
\end{proof}

\rem In fact, only right–continuity is required to represent $N$ and $D$ as cumulative
functions of positive measures. As noted in Remark \ref{rem:cont}, the total variation
function $V$ is right–continuous whenever $a$ is right–continuous. Consequently, the
Jordan decomposition still provides two non–decreasing right–continuous functions
$N$ and $D$, which can be represented as Stieltjes measures.

Therefore, the assumption that $a$ is continuous can be relaxed: it suffices to assume
that $a$ is right–continuous and of bounded variation. In this case $a$ admits a
representation as a Stieltjes integral, and in particular $a$ is a càdlàg function.

\end{document}
